% Standalone problem document, generated from the adjacent README.md.
% Compile with XeLaTeX. Mathematical content is maintained in Markdown.
\documentclass[11pt,a4paper]{article}
\usepackage[fontsize=10.5pt]{scrextend}
\usepackage[margin=22mm,headheight=15pt,headsep=9mm,footskip=11mm]{geometry}
\usepackage{fontspec}
\setmainfont{texgyrepagella}[Extension=.otf,UprightFont=*-regular,BoldFont=*-bold,ItalicFont=*-italic,BoldItalicFont=*-bolditalic]
\setsansfont{texgyreheros}[Extension=.otf,UprightFont=*-regular,BoldFont=*-bold,ItalicFont=*-italic,BoldItalicFont=*-bolditalic]
\setmonofont{texgyrecursor}[Extension=.otf,UprightFont=*-regular,BoldFont=*-bold,ItalicFont=*-italic,BoldItalicFont=*-bolditalic,Scale=MatchLowercase]
\usepackage{amsmath,amssymb}
\usepackage{unicode-math}
\setmathfont{texgyrepagella-math.otf}
\usepackage{xcolor,microtype,enumitem,needspace,fancyhdr}
\usepackage{bookmark}
\definecolor{ink}{HTML}{17344B}
\definecolor{accent}{HTML}{197D80}
\definecolor{muted}{HTML}{596773}
\hypersetup{colorlinks=true,linkcolor=accent,urlcolor=accent,pdftitle={IE-04 - Exponential
smoothed tail bounds for partial
pivoting},pdfauthor={Open Problems in Numerical Linear Algebra}}
\urlstyle{same}
\setlength{\parindent}{0pt}
\setlength{\parskip}{5pt plus 1pt minus 1pt}
\linespread{1.04}
\setlength{\emergencystretch}{3em}
\widowpenalty=10000
\clubpenalty=10000
\displaywidowpenalty=10000
\allowdisplaybreaks[1]
\setlist{nosep,leftmargin=1.5em}
\providecommand{\tightlist}{\setlength{\itemsep}{0pt}\setlength{\parskip}{0pt}}
\providecommand{\pandocbounded}[1]{#1}
\pagestyle{fancy}
\fancyhf{}
\fancyhead[L]{\sffamily\footnotesize\color{muted}OPEN PROBLEMS IN NUMERICAL LINEAR ALGEBRA}
\fancyhead[R]{\sffamily\bfseries\footnotesize\color{accent}IE-04}
\fancyfoot[L]{\parbox[b]{0.9\textwidth}{\sffamily\fontsize{8}{10}\selectfont\color{muted}Editorial ratings; a bounded search cannot certify that a problem remains unsolved.\\Verification check: 2026-09-15}}
\fancyfoot[R]{\sffamily\footnotesize\color{muted}\thepage}
\renewcommand{\headrulewidth}{0.3pt}
\renewcommand{\headrule}{\hbox to\headwidth{\color{accent}\leaders\hrule height \headrulewidth\hfill}}
\makeatletter
\renewcommand{\section}{\Needspace{5\baselineskip}\@startsection{section}{1}{0pt}{12pt plus 2pt minus 2pt}{4pt}{\sffamily\bfseries\large\color{ink}}}
\renewcommand{\subsection}{\Needspace{4\baselineskip}\@startsection{subsection}{2}{0pt}{9pt plus 1pt minus 1pt}{3pt}{\sffamily\bfseries\normalsize\color{ink}}}
\makeatother
\setcounter{secnumdepth}{0}
\begin{document}
{\sffamily\small\color{muted}Linear systems and elimination\par}
\vspace{3pt}
{\raggedright\hyphenpenalty=10000\exhyphenpenalty=10000\sffamily\bfseries\fontsize{21}{24}\selectfont\color{ink}Exponential
smoothed tail bounds for partial pivoting\par}
\vspace{7pt}
{\sffamily\small\textbf{Difficulty:} extreme\quad\textbf{Importance:} broadly
interesting\par}
{\sffamily\small\bfseries\color{accent}Status: Lean verified\par}
\vspace{4pt}
{\color{accent}\hrule height 0.8pt}
\vspace{6pt}
\textbf{Rating rationale:} Extreme reflects a uniform tail bound over
adversarial matrix centers, beyond current average-case analysis; broad
impact is justified by explaining the stability of a standard dense
solver.

\subsection{Resolution --- 11 September
2026}\label{resolution-11-september-2026}

\textbf{Solved negatively.} George Stepaniants, Department of Computing
and Mathematical Sciences, California Institute of Technology, Pasadena,
California, USA, gives a counterexample to the displayed universal
exponential tail. The
\href{https://github.com/ajt60gaibb/OpenProblemsInNLA/blob/main/linear-systems-and-elimination/IE-04/solution.md}{complete
proof's theorem, equation (2), and robustness lemma} show that, for
every \(n\ge2\),

\[\Pr\!\left\{\rho_{\mathrm{PP}}(I_n+G)>\tfrac12(3/2)^{n-1}\right\}
\ge 2^{-n^2(n^2+n+5)}.\]

For any proposed \(c_1,c_2>0\), take \(\bar A=I_n\), \(\sigma=1\) and
\(x=(3/2)^{n-1}/(2n^{c_1})\). Sufficiently large \(n\) give an
admissible \(x\ge1\) and violate the proposed upper bound. This covers
the exact real Gaussian model and every proposed universal pair; the
event has strict pivot choices and nonsingular matrices. The conclusion
uses the unrestricted range of \(x\) and does not address a different
tail restricted to smaller \(x\).

\href{https://github.com/ajt60gaibb/OpenProblemsInNLA/blob/main/linear-systems-and-elimination/IE-04/solution.pdf}{Proof
PDF} ·
\href{https://github.com/ajt60gaibb/OpenProblemsInNLA/blob/main/linear-systems-and-elimination/IE-04/solution.tex}{Standalone
source} ·
\href{https://github.com/ajt60gaibb/OpenProblemsInNLA/blob/main/references/stepaniants-ie04-2026-09-11/verification/IE-04-independent-review.md}{Independent
complete-proof review} ·
\href{https://github.com/ajt60gaibb/OpenProblemsInNLA/blob/main/references/stepaniants-ie04-2026-09-11/README.md}{Submission
and verification record}. AI assistance and automated-review limits are
disclosed. The permanent ID, original statement, and historical ratings
below are retained.

\subsection{Lean proof and verification evidence --- 15 September
2026}\label{lean-proof-and-verification-evidence-15-september-2026}

\textbf{Formalization: George Stepaniants}, Department of Computing and
Mathematical Sciences, California Institute of Technology. The
\href{https://github.com/sgstepaniants/OpenProblemsInNLA/tree/026b3e5534a4d6e15ebffb85318c2ff031df32bf/linear-systems-and-elimination/IE-04/lean}{immutable
complete proof} verifies the full negative answer through 21 checked
declarations, ending with \textit{NLA.IE04.counterexample} and
\textit{NLA.IE04.not\_uniformExponentialTail}. It covers arbitrary
positive real \(c_1,c_2\), actual Gaussian product measures, the full
nonsingular perturbation box and every admissible partial-pivoting tie
rule. The center \(I_n\) and noise scale one are admissible witnesses;
all original quantifiers are retained.

\href{https://github.com/sgstepaniants/OpenProblemsInNLA/actions/runs/35034399633/job/104600154206}{Canonical
non-root Linux run 35034399633} accepted these exact proof bytes with
Lean 4.33.1, Mathlib \textit{0df444a360eaa60ab8c11dca51a86af692955474}
and LeanCert \textit{621a43d7cf21f87872392a01e874f2f1dbddc926}. LeanCert
checks the consumed scalar certificate \(e^{-2}>1/8\) in kernel mode.
Comparator, independent default-kernel replay and the required rejection
and sandbox controls passed; only \textit{propext},
\textit{Classical.choice} and \textit{Quot.sound} occur in the exported
proofs.

\href{https://github.com/ajt60gaibb/OpenProblemsInNLA/blob/main/linear-systems-and-elimination/IE-04/lean/README.md}{Project
and reproduction instructions} ·
\href{https://github.com/ajt60gaibb/OpenProblemsInNLA/blob/main/linear-systems-and-elimination/IE-04/lean/verification/linux-2026-09-15}{Original
runtime evidence} ·
\href{https://github.com/ajt60gaibb/OpenProblemsInNLA/blob/main/linear-systems-and-elimination/IE-04/lean/reviews/final}{Both
independent final referee reports} ·
\href{https://github.com/ajt60gaibb/OpenProblemsInNLA/blob/main/linear-systems-and-elimination/IE-04/lean/formalization.yaml}{Formalization
manifest}. Substantial AI assistance and scoped AI-agent review are
disclosed. The checked proof revision is distinct from later
documentation commits, whose publication workflow must run separately.
No external human peer review or official Tau Ceti endorsement is
asserted.

\newpage

\subsection{Context and notation}\label{context-and-notation}

All elimination in this problem is in exact arithmetic. Write
\(\|A\|_{\max}=\max_{ij}|a_{ij}|\). A pivoting path creates successive
active Schur complements \(S_1=A,S_2,\ldots,S_n\), with row/column
permutations as appropriate. Its element-growth factor is

\[\rho(A)=\frac{\max_{1\leq j\leq n}\|S_j\|_{\max}}{\|A\|_{\max}}.\]

Partial pivoting chooses a largest-magnitude entry in the active first
column; complete pivoting chooses one anywhere in the active matrix. If
ties occur, a universal statement includes every admissible tie choice;
a supremum includes all admissible paths. These conventions remove
implementation-dependent ambiguity. Matrices are nonsingular unless
otherwise stated.

\subsection{Problem statement}\label{problem-statement}

Do universal constants \(c_1,c_2>0\) exist such that, for every
\(n\geq1\), deterministic \(\bar A\in\mathbb R^{n\times n}\) with
\(\|\bar A\|_2\leq1\), \(0<\sigma\leq1\), and \(x\geq1\), the matrix
\(A=\bar A+\sigma G\), with independent standard normal entries in
\(G\), satisfies

\[\Pr\!\left\{\rho_{\mathrm{PP}}(A)>x(n/\sigma)^{c_1}\right\}
\leq 2^{-c_2x}?\]

This is a uniform smoothed-analysis question: the deterministic center
may itself be a worst-case input. Almost surely the perturbed matrix is
nonsingular and the pivot choices have no ties. Numerical evidence or
estimates only at \(\bar A=0\) do not settle the quantifier over
deterministic centers. The tail estimate would quantify the rarity of
substantial growth after small Gaussian input perturbations.

\subsection{References}\label{references}

Spielman and Teng,
\href{https://www.cs.yale.edu/homes/spielman/PAPERS/focmSmoothed.pdf}{\emph{Smoothed
Analysis of Algorithms and Heuristics: Progress and Open Questions}},
§P6, Conjecture 16, p.~52 of the author PDF; in \emph{Foundations of
Computational Mathematics, Santander 2005} (2006), 274--342,
\href{https://doi.org/10.1017/CBO9780511721571.010}{chapter DOI}. Huang
and Tikhomirov,
\href{https://doi.org/10.1007/s00440-024-01276-2}{\emph{Average-case
analysis of the Gaussian elimination with partial pivoting}}, PTRF 189
(2024), introduction.

\subsection{Earlier status check ---
2026-09-08}\label{earlier-status-check-2026-09-08}

Searches for \textit{Exponential Stability of GEPP solved} and
\textit{Gaussian elimination smoothed analysis 2025 2026} found the
average-case theorem, not this exponential tail bound. The 2026
Peca-Medlin butterfly paper also distinguishes average-case results from
the still-unavailable full smoothed analysis. Randomizing the pivot rule
is a different model.

\subsection{Audit update --- 2026-09-10}\label{audit-update-2026-09-10}

The linked Spielman--Teng chapter was checked against the
\href{https://www.cs.yale.edu/homes/spielman/SmoothedAnalysis/surveys.html}{author
bibliography}: its title has been corrected in the reference, with
Conjecture 16 unchanged.
\href{https://link.springer.com/article/10.1007/s00440-024-01276-2}{Huang--Tikhomirov}
explicitly discusses the obstruction to extending its Gaussian result to
arbitrary centers. Searches for exponential GEPP tails and subsequent
smoothed-analysis results found no theorem with the displayed uniform
quantifiers.
\end{document}
